\section*{Capítulo \ref{ch:b2:neurons-synapses} --- Neuronas, potenciales de acción y sinapsis}

\begin{solution}{exo:b2:neurons-synapses:1}
Las \omterm{def:b2:neurons-synapses:neuron}{dendritas} reciben las \omterm{def:b2:neurons-synapses:synapse}{sinapsis}; el soma contiene el núcleo e integra;
el \omterm{prop:b2:neurons-synapses:integration}{cono axónico} decide; el \omterm{def:b2:neurons-synapses:neuron}{axón} conduce; las terminales liberan el
transmisor. La \omterm{def:b2:neurons-synapses:neuron}{mielina} es la membrana enrollada de una célula glial, casi
sin citoplasma; aísla el \omterm{def:b2:neurons-synapses:neuron}{axón} de modo que el impulso salta de un nódulo a
otro y viaja cien veces más deprisa con un coste menor.
\end{solution}

\begin{solution}{exo:b2:neurons-synapses:2}
Despolarización hasta el umbral; fase ascendente --- se abren los canales
de sodio dependientes de voltaje, retroalimentación positiva; máximo cerca
de $E_{\mathrm{Na}}$ --- los canales de sodio se inactivan; fase
descendente --- se abren los canales de potasio dependientes de voltaje;
hiperpolarización posterior --- los canales de potasio siguen abiertos,
cerca de
$E_{\mathrm K}$; recuperación a medida que se cierran. De todo o nada
porque la retroalimentación positiva llega hasta el final o no se produce;
en un solo sentido porque la membrana recién recorrida está en \omterm{prop:b2:neurons-synapses:ap}{periodo
refractario} (canales de sodio inactivados).
\end{solution}

\begin{solution}{exo:b2:neurons-synapses:3}
El impulso llega a la terminal; se abren los canales de calcio dependientes
de voltaje; el calcio desencadena la fusión de las vesículas; el
transmisor difunde a través de la hendidura; se une a los receptores
postsinápticos; se abren canales (o se activa una \omterm{prop:b2:cell-signalling:gpcr}{proteína G}); \omterm{def:b2:neurons-synapses:synapse}{potencial
postsináptico}; el transmisor se retira mediante una enzima o un
transportador; se recupera la membrana de las vesículas.
\end{solution}

\begin{solution}{exo:b2:neurons-synapses:4}
Eléctrica: sin retraso, normalmente en ambos sentidos, sin amplificación,
sin inhibición (la corriente solo pasa). Química: medio milisegundo, en un
solo sentido, con amplificación (un impulso libera cientos de cuantos),
inhibición posible (canales de cloruro o de potasio), y modificable.
\end{solution}

\begin{solution}{exo:b2:neurons-synapses:5}
A partir de $E = (61.5/z)\log_{10}(c_{\text{ext}}/c_{\text{int}})$:
potasio, \qty{-89}{mV}; sodio, \qty{+61}{mV}; cloruro, con $z =
-1$, $-61.5\log_{10} 11 = \qty{-64}{mV}$; calcio, con $z = 2$,
$30.75\times 4.3 = \qty{+132}{mV}$. A \qty{-70}{mV}: el potasio sale, el
sodio entra, el cloruro sale ligeramente (la membrana está por debajo de
su \omterm{thm:b2:neurons-synapses:nernst}{potencial de equilibrio}) y el calcio entra con fuerza.
\end{solution}

\begin{solution}{exo:b2:neurons-synapses:6}
$20 : 1$: $(20\times(-89) + 61)/21 = \qty{-82}{mV}$, el estado de reposo,
cerca de $E_{\mathrm K}$. $1 : 20$: $(-89 + 20\times 61)/21 =
\qty{+54}{mV}$, el máximo, cerca de $E_{\mathrm{Na}}$. $50 : 1$: $(50\times
(-89) + 61)/51 = \qty{-86}{mV}$, la hiperpolarización posterior, con
canales de potasio adicionales abiertos.
\end{solution}

\begin{solution}{exo:b2:neurons-synapses:7}
$\sqrt{500} = 22$: \qty{22}{m/s}. Para alcanzar \qty{100}{m/s}, un \omterm{def:b2:neurons-synapses:neuron}{axón}
amielínico necesitaría $d = 10^{4}$ \unit{\micro m}, un
centímetro. Un metro lleva \qty{10}{ms} en la fibra mielínica,
\qty{45}{ms} en el \omterm{def:b2:neurons-synapses:neuron}{axón} de calamar y un segundo entero en una fibra de
\qty{1}{\micro m}.
\end{solution}

\begin{solution}{exo:b2:neurons-synapses:8}
$m = \ln(200/74) = 1.0$; $P(1) = e^{-1} = 0.37$, $P(2) = 0.18$, $P(3)
= 0.06$. Respuesta media: $1.0\times 0.5 = \qty{0.5}{mV}$.
\end{solution}

\begin{solution}{exo:b2:neurons-synapses:9}
$1/\qty{2}{ms} = 500$ impulsos por segundo. La intensidad del contacto se
codifica mediante la frecuencia de los impulsos (y mediante el número de
\omterm{def:b2:neurons-synapses:neuron}{neuronas} reclutadas); los impulsos en sí son idénticos.
\end{solution}

\begin{solution}{exo:b2:neurons-synapses:10}
$\lambda = \sqrt{R_m d/4R_i}$: \qty{0.5}{mm}, \qty{1.6}{mm},
\qty{11}{mm}. Con \omterm{def:b2:neurons-synapses:neuron}{mielina}, a \qty{10}{\micro m}: $1.6\times\sqrt{200} =
\qty{22}{mm}$. Un internodo de \qty{1}{mm} solo pierde $1 - e^{-1/22}
= \qty{4}{\%}$ de la señal, de modo que el nódulo siguiente llega
fácilmente al umbral. A través de \qty{3}{mm} desmielinizados, la señal
cae a $e^{-3/1.6} = \qty{15}{\%}$ --- de un impulso de \qty{100}{mV},
aproximadamente el umbral de \qty{15}{mV}: la conducción falla o se vuelve
poco fiable.
\end{solution}

\begin{solution}{exo:b2:neurons-synapses:11}
Tetrodotoxina: sin corriente de sodio, sin impulso, sin transmisión.
Tetraetilamonio: sin corriente de potasio, impulso prolongado, más
transmisor liberado por impulso. Sin calcio: impulsos normales, sin
liberación, sin transmisión. Curare: liberación normal, receptores
bloqueados, sin potencial de placa terminal --- parálisis. Inhibidor de la
colinesterasa: la acetilcolina persiste, los receptores siguen abiertos,
el músculo se despolariza y luego no puede repolarizarse --- parálisis
por exceso. Toxina botulínica: las vesículas no pueden fusionarse, no hay
liberación --- parálisis.
\end{solution}

\begin{solution}{exo:b2:neurons-synapses:12}
El retraso compra: la transmisión en un solo sentido; la amplificación
(un solo impulso libera cientos de cuantos y puede activar una célula
mayor); la inhibición, que una unión eléctrica no puede producir; y la
plasticidad, puesto que el número de cuantos y de receptores puede
cambiar con el uso, que es como aprenden los circuitos. Las \omterm{def:b2:neurons-synapses:synapse}{sinapsis}
eléctricas se usan donde la sincronía y la velocidad importan más que el
cálculo: circuitos de huida, el músculo cardíaco, algunas redes
inhibidoras.
\end{solution}

\begin{solution}{pb:b2:neurons-synapses:1}
\textbf{1.} $E_{\mathrm K} = 61.5\log_{10}(4/140) = \qty{-95}{mV}$;
$E_{\mathrm{Na}} = 61.5\log_{10}(145/12) = \qty{+67}{mV}$.
\textbf{2.} $(25\times(-95) + 67)/26 = \qty{-89}{mV}$.
\textbf{3.} $E_{\mathrm K} = 61.5\log_{10}(8/140) = \qty{-76}{mV}$;
$V = (25\times(-76) + 67)/26 = \qty{-71}{mV}$: más cerca del umbral, de
modo que al principio es más excitable --- y, si persiste, menos, porque
los canales de sodio se inactivan en el nivel despolarizado.
\textbf{4.} Los gradientes se disipan a lo largo de horas a medida que el
sodio entra y el potasio sale; el potencial deriva hacia cero; con la
bomba detenida, los aniones no difusibles de la célula atraen cloruro y
agua, y la célula se hincha.
\textbf{5.} Superficie por milímetro: $\pi\times 15\times 10^{-6}\times
10^{-3} = \qty{4.7e-8}{m^2}$; $C = \qty{4.7e-10}{F}$; $Q = CV =
\qty{4.7e-11}{C}$: $2.9\times 10^{8}$ iones.
\textbf{6.} Volumen por milímetro: $\pi(7.5\times 10^{-6})^{2}\times
10^{-3} = \qty{1.8e-13}{m^3} = \qty{1.8e-10}{L}$, que contiene $12\times
10^{-3}\times 1.8\times 10^{-10} = 2.1\times 10^{-12}$ mol, $1.3\times
10^{12}$ iones: un cambio de $2\times 10^{-4}$.
\textbf{7.} $2.9\times 10^{8}/3 \approx 10^{8}$ ATP por milímetro y por
impulso.
\textbf{8.} $0.9/90 = \qty{10}{ms}$ en cada sentido.
\textbf{9.} Desnudo: $\sqrt{1\times 15\times 10^{-6}/4} = \qty{1.9}{mm}$;
con \omterm{def:b2:neurons-synapses:neuron}{mielina}: $\times\sqrt{250} = \qty{31}{mm}$.
\textbf{10.} Desnudo: $e^{-1.5/1.9} = 0.45$, \qty{45}{mV}; con \omterm{def:b2:neurons-synapses:neuron}{mielina}:
$e^{-1.5/31} = 0.95$, \qty{95}{mV}. Ambos superan el umbral de
\qty{15}{mV}; pero el \omterm{def:b2:neurons-synapses:neuron}{axón} desnudo debe cargar toda su membrana por el
camino, que es lo que lo hace lento, mientras que bajo la \omterm{def:b2:neurons-synapses:neuron}{mielina} casi
nada se pierde ni se carga.
\textbf{11.} $\sqrt{15} = \qty{3.9}{m/s}$: \qty{0.23}{s} en cada sentido,
casi medio segundo para la ida y la vuelta --- demasiado lento para un
reflejo que debe corregir un tropiezo.
\textbf{12.} $e^{-4/1.9} = 0.12$: \qty{12}{mV} en el nódulo lejano, por
debajo del umbral --- el impulso queda bloqueado y el reflejo se pierde.
\textbf{13.} Detrás del impulso, los canales de sodio quedan inactivados
durante uno o dos milisegundos, de modo que la membrana recién recorrida no
puede volver a excitarse y la onda solo puede avanzar.
\textbf{14.} $m = \ln(300/111) = 1.0$.
\textbf{15.} $P(1) = 0.37$, $P(2) = 0.18$, $P(3) = 0.06$: unos 110, 55
y 18 de los 300 ensayos.
\textbf{16.} $1.0\times 0.4 = \qty{0.4}{mV}$.
\textbf{17.} $P(0) = e^{-60} \approx 10^{-26}$: nunca; respuesta media de
\qty{24}{mV}, por encima del umbral de \qty{15}{mV}: una sola de estas
\omterm{def:b2:neurons-synapses:synapse}{sinapsis} hace disparar la motoneurona.
\textbf{18.} Veinte \omterm{def:b2:neurons-synapses:synapse}{sinapsis} activas a la vez suman sus corrientes
(sumación espacial), de forma sublineal a medida que el potencial se
acerca al potencial de inversión de los canales; incluso con $m = 1$ cada
una darían \qty{8}{mV}, y con la $m$ normal mucho más que el umbral.
\textbf{19.} Abrir canales de cloruro con $E_{\mathrm{Cl}} = V_{\text{rep}}$
añade conductancia sin mover el potencial; la corriente excitadora fluye
ahora a través de una membrana con más fugas y produce una
despolarización menor --- inhibición por cortocircuito.
\textbf{20.} $10 + 0.6 + 10 + 0.6 + 5 = \qty{26}{ms}$, frente a los
\qty{30}{ms} medidos.
\textbf{21.} La propia activación del receptor en el huso muscular, la
propagación del \omterm{prop:b2:neurons-synapses:ap}{potencial de acción} de la \omterm{def:b2:cell-differentiation:myogenesis}{fibra muscular} a lo largo de la
fibra (metros por segundo sobre centímetros) y la latencia mecánica antes
de que aparezca la fuerza.
\textbf{22.} Cada \omterm{def:b2:neurons-synapses:synapse}{sinapsis} añade medio milisegundo, y una interneurona
añade una célula; pero las interneuronas permiten invertir la señal
(inhibiendo el antagonista), combinarla con otras y que el encéfalo la
module --- un reflejo que puede moldearse.
\textbf{23.} $90 - 6.2 = \qty{84}{ms}$ de conducción para \qty{1.8}{m}:
unos \qty{21}{m/s}.
\textbf{24.} La mitad de la corriente entrante a cada voltaje: el umbral
sube, el impulso es más pequeño y más lento, el margen de seguridad en cada
nódulo disminuye y el impulso puede fallar; el reflejo se debilita o
desaparece.
\textbf{25.} \omterm{thm:b2:neurons-synapses:chord}{Potencial de reposo}, \qty{-89}{mV}; \qty{10}{ms} en cada
sentido; $m = 1.0$; latencia calculada, \qty{26}{ms}.
\end{solution}
